\chapter{Hipótesis y DSL}
\label{sec:man-dsl}

\section{La idea}

Una \emph{hipótesis} es una pregunta sobre el grafo expresada en
términos de entidades y relaciones. El DSL describe ``formas''
que el sistema busca:

\begin{lstlisting}[basicstyle=\ttfamily\small]
Entity -[Relation]-> Entity
\end{lstlisting}

Leída como ``una \code{Entity} conectada por una \code{Relation} a
otra \code{Entity}''. El resultado son todos los paths concretos
del grafo que encajan.

\section{Sintaxis mínima}

\subsection{Un solo paso}

\begin{lstlisting}[basicstyle=\ttfamily\small]
User -[Auth]-> Host
\end{lstlisting}

Busca usuarios que se autenticaron en algún host.

\subsection{Cadena de pasos}

\begin{lstlisting}[basicstyle=\ttfamily\small]
User -[Auth]-> Host -[Connect]-> Host {k=3}
\end{lstlisting}

Usuario que autenticó contra un host, y desde ese host se abrió
una conexión a otro. \code{k=3} es una cota de simplicidad: acota
la explosión combinatoria limitando longitud o ramificación.

\subsection{Comodín}

\begin{lstlisting}[basicstyle=\ttfamily\small]
* -[Execute]-> Process
\end{lstlisting}

Cualquier entidad que ejecutó un proceso. Útil cuando no sabés de
dónde mirar.

\subsection{Predicados}

\begin{lstlisting}[basicstyle=\ttfamily\small]
Process {name="cmd.exe"} -[Spawn]-> Process {name="powershell.exe"}
\end{lstlisting}

Restringe los nodos por propiedad. Comparadores: \code{=},
\code{!=}, \code{>}, \code{<}, \code{>=}, \code{<=},
\code{like} (strings).

\subsection{Ventana temporal}

\begin{lstlisting}[basicstyle=\ttfamily\small]
Host -[DNS]-> Domain {matches_regex(value, "^[a-f0-9]{32}\.")}
{window=1h}
\end{lstlisting}

Limita la separación máxima entre eventos del path. Unidades:
\code{s}, \code{m}, \code{h}, \code{d}.

\section{Catálogo built-in}

Tipos de entidad: \code{IP}, \code{Host}, \code{User},
\code{Process}, \code{File}, \code{Domain}, \code{Registry},
\code{URL}, \code{Service}, más el comodín \code{*}.

Relaciones: \code{Auth}, \code{Connect}, \code{Execute},
\code{Read}, \code{Write}, \code{DNS}, \code{Modify},
\code{Spawn}, \code{Delete}, \code{*}.

Predicados: \code{in_subnet(ip, cidr)},
\code{matches_regex(field, pattern)}, \code{is_weekend()},
\code{within_hours(start, end)}, \code{tag(name)}.

Para la gramática formal completa ver el libro de Especificación,
cap.\ \emph{DSL}.

\section{Ejemplos clásicos}

\subsection{Lateral movement}

\begin{lstlisting}[basicstyle=\ttfamily\small]
User -[Auth]-> Host -[Connect]-> Host
  {window=30m, k=3}
\end{lstlisting}

\subsection{Exfiltración sospechosa}

\begin{lstlisting}[basicstyle=\ttfamily\small]
Process -[Read]-> File
  -[Connect]-> IP {not in_subnet(value, "10.0.0.0/8")}
  {window=1h}
\end{lstlisting}

Proceso que lee un archivo y después se conecta a una IP externa
dentro de una hora.

\subsection{Cadena de spawning Office}

\begin{lstlisting}[basicstyle=\ttfamily\small]
Process {name="winword.exe"}
  -[Spawn]-> Process {name="cmd.exe"}
  -[Spawn]-> Process {name="powershell.exe"}
\end{lstlisting}

\subsection{DNS tunneling}

\begin{lstlisting}[basicstyle=\ttfamily\small]
Host -[DNS]-> Domain
  {matches_regex(value, "^[a-f0-9]{32}\.")}
  {window=1h}
\end{lstlisting}

\section{Buenas y malas hipótesis}

\textbf{Buenas}:

\begin{itemize}
    \item Específicas sobre tipos y relaciones (no comodín en
    todo).
    \item Con \code{k} y \code{window} acotadas.
    \item Con predicados que descartan casos comunes (ej.\ excluir
    subnet corporativa para mirar tráfico externo).
\end{itemize}

\textbf{Malas}:

\begin{itemize}
    \item \code{* -[*]-> *} — devuelve todo el grafo, no informa.
    \item Sin \code{k} ni \code{window} sobre un grafo grande —
    puede tomar minutos o fallar por presión de memoria.
    \item Demasiado restrictivas — cero resultados; considerá
    relajar un predicado y revisar.
\end{itemize}

\section{Validar sin ejecutar}

En la UI, botón \emph{Parse} al lado de \emph{Run}. El servidor
valida la sintaxis y devuelve el AST si es válida, con la línea y
columna del error si no. Útil para hipótesis largas antes de
comprometerte a correrlas.

\section{Guardar hipótesis}

Una hipótesis útil se puede convertir en regla Sigma para
portarla a otros sistemas (\code{Hunt \textrightarrow{} Save as
Sigma}). También podés guardarla en la librería de hunts compartidos
de la sesión para reuso.
